\chapter{Avaliação das Instalações Elétricas}

A avaliação das instalações elétricas tem por finalidade analisar, de forma
técnica e qualitativa, as condições gerais do sistema elétrico da unidade
consumidora \textbf{<<nomeUc>>}, com foco nos \ac{QGBTs} e nos \ac{QDs}. Quando aplicável, a análise contempla também a subestação interna da unidade.

Esta avaliação busca identificar condições que possam impactar a segurança, a
confiabilidade operacional e a implantação de ações de eficiência energética.

\section{Procedimentos de Inspeção e Referências Normativas}

A avaliação das instalações elétricas foi realizada por meio de inspeção técnica
presencial, com base em observação visual, registros fotográficos e verificação
das condições aparentes dos quadros elétricos. Durante a inspeção, foram
analisados aspectos relacionados à integridade física dos painéis, organização
interna dos condutores, identificação dos circuitos, dispositivos de proteção e
condições aparentes do sistema de aterramento e equipotencialização.

As análises foram conduzidas com base nas boas práticas aplicáveis às instalações
elétricas de baixa tensão, conforme estabelecido pela NBR~5410~\cite{nbr5410}. 
Para unidades que possuem subestação própria, também foram considerados os 
requisitos aplicáveis às instalações elétricas de média tensão previstos na 
NBR~14039~\cite{nbr14039}. Adicionalmente, foram considerados os princípios de segurança aplicáveis às instalações e serviços em eletricidade, conforme a NR~10~\cite{nr10}.

\section{Avaliação dos Quadros Elétricos}


A avaliação dos quadros elétricos foi realizada por meio de inspeção visual dos 
painéis da unidade consumidora, considerando aspectos construtivos, operacionais 
e de segurança. A análise foi conduzida com base nos requisitos aplicáveis às 
instalações elétricas de baixa tensão estabelecidos pela NBR~5410~\cite{nbr5410}, 
com foco na verificação das condições de proteção elétrica, identificação dos 
circuitos, organização interna dos condutores, integridade dos invólucros, 
condições aparentes do sistema de aterramento e acessibilidade para operação e 
manutenção.

As Figuras~\ref{nr01acft}, \ref{nr02acft} e \ref{nr03acft} apresentam exemplos 
de quadros elétricos vistoriados durante a inspeção técnica, incluindo registros 
das condições externas e internas dos painéis avaliados.

\vspace{1cm}
% --- Fotos ---
\begin{figure}[H]
\centering
\includegraphics[width=0.8\textwidth]{Figuras/NR01ACFT.png}
\caption{Vista externa e interna do Quadro Elétrico Nr. 01}
\label{nr01acft}
\end{figure}

\begin{figure}[H]
\centering
\includegraphics[width=0.6\textwidth]{Figuras/NR02ACFT.png}
\caption{Vista externa e interna do Quadro Elétrico Nr. 02}
\label{nr02acft}
\end{figure}

\begin{figure}[H]
\centering
\includegraphics[width=0.6\textwidth]{Figuras/NR03ACFT.png}
\caption{Vista externa e interna do Quadro Elétrico Nr. 03}
\label{nr03acft}
\end{figure}

Após a apresentação dos registros fotográficos, o quadro a seguir
resume os principais itens verificados e as respectivas conformidades observadas 
nos painéis avaliados.

\begin{comment}
  A avaliação dos quadros elétricos foi realizada por meio de inspeção visual dos
painéis, considerando suas condições construtivas e operacionais. Para cada
quadro avaliado, foram identificadas eventuais não conformidades e apontados os
itens passíveis de correção, visando à adequação das instalações aos requisitos
aplicáveis da NBR~5410~\cite{nbr5410}.

Os comentários técnicos contemplaram a verificação dos seguintes aspectos:

\begin{itemize}
    \item Atendimento aos requisitos mínimos de segurança elétrica;
    \item Existência de dispositivos de proteção contra surtos;
    \item Identificação dos circuitos e disponibilidade de diagrama elétrico;
    \item Organização interna dos condutores e adequação do invólucro;
    \item Condições aparentes do sistema de aterramento e equipotencialização;
    \item Condições de acesso para operação e manutenção.
\end{itemize}

As não conformidades identificadas foram registradas de forma objetiva,
acompanhadas de recomendações técnicas, com foco na melhoria da segurança, da
confiabilidade operacional e da adequação das instalações elétricas. A avaliação da subestação foi conduzida com base nos requisitos aplicáveis às instalações elétricas de média tensão, conforme a NBR~14039~\cite{nbr14039}.  
\end{comment}




\begin{comment}

\subsection{\textbf{\textcolor{red}{EXEMPLO - Quadro Elétrico: ESCOLA001 – Nr. 01}} }

Os dados a seguir referem-se ao quadro elétrico \textbf{\textcolor{red}{Identificação do quadro (EX: Nr. 01)}}, avaliado por meio de inspeção visual, conforme os critérios estabelecidos neste
capítulo.

% --- Tabela ---
\begin{table}[H]
    \centering
    \includegraphics[width=0.8\textwidth]{Figuras/NR01AC.png}
    \caption{Dados técnicos — Quadro Elétrico Nr. 01}
    \label{nr01ac}
\end{table}

% --- Fotos ---
\begin{figure}[H]
\centering
\includegraphics[width=0.8\textwidth]{Figuras/NR01ACFT.png}
\caption{Vista externa e interna do Quadro Elétrico Nr. 01}
\label{nr01acft}
\end{figure}

\subsubsection{Comentários Técnicos}

\textcolor{red}{Nesta subseção devem ser descritas as principais condições observadas no quadro elétrico avaliado, destacando eventuais não conformidades em relação aos requisitos estabelecidos pela NBR~5410~\cite{nbr5410}. Devem ser abordados, de forma objetiva, aspectos como a existência ou não de dispositivos de proteção contra surtos (\ac{DPS}), identificação dos circuitos, organização interna dos condutores, adequação do invólucro, condições aparentes do sistema de aterramento e facilidades para operação e manutenção.}

\textcolor{red}{Quando aplicável, podem ser mencionadas observações adicionais relevantes para a segurança e a confiabilidade operacional das instalações, evitando a inclusão de valores ou medições detalhadas que não façam parte do escopo do diagnóstico.}

\subsubsection{Recomendações}

\textcolor{red}{Nesta subseção devem ser apresentadas recomendações técnicas decorrentes das não conformidades identificadas, com foco na adequação do quadro elétrico aos requisitos da NBR~5410~\cite{nbr5410}. As recomendações devem ser claras, objetivas e compatíveis com o escopo do diagnóstico energético, podendo incluir, por exemplo, a necessidade de instalação de \ac{DPS}, reorganização interna dos condutores, atualização da identificação dos circuitos, adequação do invólucro ou substituição do painel, quando tecnicamente justificável.}

\end{comment}





{\footnotesize
\renewcommand{\arraystretch}{1.25}
\arrayrulecolor{gridcol}
\noindent
\resizebox{\textwidth}{!}{%
\begin{tabular}{|C{2.3cm}|L{7cm}|C{1.6cm}|C{2cm}|C{2cm}|}
\hline
\rowcolor{headbg}
\textcolor{white}{\textbf{Ref. NBR 5410:2004}} &
\textcolor{white}{\textbf{Item de Verificação}} &
\textcolor{white}{\textbf{Conforme}} &
\textcolor{white}{\textbf{Não Conforme}} &
\textcolor{white}{\textbf{Não Aplicável}} \\ \hline

\rowcolor{secbg}
\multicolumn{5}{|l|}{\textbf{SEGURANÇA ELÉTRICA E PROTEÇÕES}} \\ \hline
6.3.5.2.1 & Possui Dispositivo de Proteção contra Surto (DPS) &
\cellcolor{okbg} & \cellcolor{nokbg} & \cellcolor{nabg} \\ \hline
8.3.2.2 & Disjuntores em bom estado de funcionamento &
\cellcolor{okbg} & \cellcolor{nokbg} & \cellcolor{nabg} \\ \hline
5.3.4.1 & Disjuntores corretamente dimensionados (sem sinais de sobrecarga) &
\cellcolor{okbg} & \cellcolor{nokbg} & \cellcolor{nabg} \\ \hline

\rowcolor{secbg}
\multicolumn{5}{|l|}{\textbf{IDENTIFICAÇÃO E DOCUMENTAÇÃO}} \\ \hline
6.1.8.1 & Diagrama unifilar disponível e atualizado no quadro &
\cellcolor{okbg} & \cellcolor{nokbg} & \cellcolor{nabg} \\ \hline
6.5.4.9 & Circuitos identificados de forma clara e legível, com etiquetas em bom estado &
\cellcolor{okbg} & \cellcolor{nokbg} & \cellcolor{nabg} \\ \hline
6.5.4.11 & Plaqueta de advertência de risco elétrico fixada &
\cellcolor{okbg} & \cellcolor{nokbg} & \cellcolor{nabg} \\ \hline

\rowcolor{secbg}
\multicolumn{5}{|l|}{\textbf{INVÓLUCRO E ORGANIZAÇÃO INTERNA}} \\ \hline
8.3.2.1 & Invólucro (caixa/gabinete) em bom estado de conservação &
\cellcolor{okbg} & \cellcolor{nokbg} & \cellcolor{nabg} \\ \hline
5.1.2.2.2 & Proteção de policarbonato cobrindo todas as partes energizadas &
\cellcolor{okbg} & \cellcolor{nokbg} & \cellcolor{nabg} \\ \hline
6.2.9.6.2 & Eletrodutos vedados com espuma expansiva antichamas &
\cellcolor{okbg} & \cellcolor{nokbg} & \cellcolor{nabg} \\ \hline
8.3.1 & Condutores organizados, sem excesso de fiação solta &
\cellcolor{okbg} & \cellcolor{nokbg} & \cellcolor{nabg} \\ \hline
6.1.5.3 & Condutores devidamente identificados (cores normativas NBR 5410) &
\cellcolor{okbg} & \cellcolor{nokbg} & \cellcolor{nabg} \\ \hline
8.3.1 & Não há condutores com isolação danificada ou ressecada &
\cellcolor{okbg} & \cellcolor{nokbg} & \cellcolor{nabg} \\ \hline
6.1.4 & Espaço livre adequado para operação e manutenção &
\cellcolor{okbg} & \cellcolor{nokbg} & \cellcolor{nabg} \\ \hline
8.3.2.1 & Não há materiais estranhos armazenados no interior &
\cellcolor{okbg} & \cellcolor{nokbg} & \cellcolor{nabg} \\ \hline

\rowcolor{secbg}
\multicolumn{5}{|l|}{\textbf{ATERRAMENTO E EQUIPOTENCIALIZAÇÃO}} \\ \hline
6.4.2.1 & Barramento PE presente e identificado &
\cellcolor{okbg} & \cellcolor{nokbg} & \cellcolor{nabg} \\ \hline
8.3.1 & Conexões de aterramento sem oxidação ou folga &
\cellcolor{okbg} & \cellcolor{nokbg} & \cellcolor{nabg} \\ \hline
6.4.4 & Equipotencialização das estruturas metálicas &
\cellcolor{okbg} & \cellcolor{nokbg} & \cellcolor{nabg} \\ \hline

\rowcolor{secbg}
\multicolumn{5}{|l|}{\textbf{CONDIÇÕES DE ACESSO E MANUTENÇÃO}} \\ \hline
6.5.4.8 & Quadro de fácil acesso (sem obstruções) &
\cellcolor{okbg} & \cellcolor{nokbg} & \cellcolor{nabg} \\ \hline
8.3.2.1 & Portas e fechamento em bom estado &
\cellcolor{okbg} & \cellcolor{nokbg} & \cellcolor{nabg} \\ \hline
6.1.4 & Iluminação adequada na área do quadro &
\cellcolor{okbg} & \cellcolor{nokbg} & \cellcolor{nabg} \\ \hline
Tabela 59 & Espaço para instalação de módulo IoT / telemetria &
\cellcolor{okbg} & \cellcolor{nokbg} & \cellcolor{nabg} \\ \hline

\end{tabular}%
}}



\section{Avaliação da Subestação}

Para unidades que possuem subestação própria, foi realizada inspeção visual das
instalações de média tensão, com o objetivo de avaliar suas condições gerais de
segurança, conservação e operação.

A análise foi conduzida com base nos requisitos aplicáveis às instalações
elétricas de média tensão estabelecidos pela NBR~14039~\cite{nbr14039}, com
foco na verificação das condições do padrão de entrada, transformador,
dispositivos de proteção e seccionamento, sistema de aterramento e condições de
acesso para operação e manutenção.

As Figuras~\ref{sub01}, \ref{sub02} e \ref{sub03} apresentam exemplos de
elementos vistoriados durante a inspeção técnica da subestação.

\vspace{1cm}

% --- Fotos ---
\begin{figure}[H]
\centering
\includegraphics[width=0.75\textwidth]{Figuras/sub01.png}
\caption{Vista geral da subestação}
\label{sub01}
\end{figure}

\begin{figure}[H]
\centering
\includegraphics[width=0.6\textwidth]{Figuras/sub02.png}
\caption{Transformador da unidade}
\label{sub02}
\end{figure}

\begin{figure}[H]
\centering
\includegraphics[width=0.6\textwidth]{Figuras/sub03.png}
\caption{Dispositivos de proteção e acesso à subestação}
\label{sub03}
\end{figure}

Após a apresentação dos registros fotográficos, o quadro a seguir
resume os principais itens verificados e as respectivas conformidades observadas
na subestação avaliada.



{\footnotesize
\renewcommand{\arraystretch}{1.2}
\arrayrulecolor{gridcol}
\noindent
\resizebox{\textwidth}{!}{%
\begin{tabular}{|C{2.3cm}|L{7cm}|C{1.6cm}|C{2cm}|C{2cm}|}
\hline
\rowcolor{headbg}
\textcolor{white}{\textbf{Ref. NBR 14039:2021}} &
\textcolor{white}{\textbf{Item de Verificação}} &
\textcolor{white}{\textbf{Conforme}} &
\textcolor{white}{\textbf{Não Conforme}} &
\textcolor{white}{\textbf{Não Aplicável}} \\ \hline

\rowcolor{secbg}
\multicolumn{5}{|l|}{\textbf{PADRÃO DE ENTRADA}} \\ \hline
8.2.2 & Padrão de entrada em bom estado de conservação &
\cellcolor{okbg} & \cellcolor{nokbg} & \cellcolor{nabg} \\ \hline
9.1.6 & Padrão protegido fisicamente (gradil, cadeados em bom estado) &
\cellcolor{okbg} & \cellcolor{nokbg} & \cellcolor{nabg} \\ \hline
5.2.2 & Eletrodutos vedados com espuma expansiva antichamas &
\cellcolor{okbg} & \cellcolor{nokbg} & \cellcolor{nabg} \\ \hline
--- & Medidor em bom estado e acessível à distribuidora &
\cellcolor{okbg} & \cellcolor{nokbg} & \cellcolor{nabg} \\ \hline
--- & Ramal de ligação em bom estado &
\cellcolor{okbg} & \cellcolor{nokbg} & \cellcolor{nabg} \\ \hline

\rowcolor{secbg}
\multicolumn{5}{|l|}{\textbf{TRANSFORMADOR}} \\ \hline
8.2.2 & Transformador em bom estado aparente (sem avarias ou vazamentos) &
\cellcolor{okbg} & \cellcolor{nokbg} & \cellcolor{nabg} \\ \hline
6.5.1 & Tipo e potência adequados à carga instalada &
\cellcolor{okbg} & \cellcolor{nokbg} & \cellcolor{nabg} \\ \hline
6.1.5 & Plaqueta de identificação legível &
\cellcolor{okbg} & \cellcolor{nokbg} & \cellcolor{nabg} \\ \hline
8.2.2 & Sem ruídos anormais ou vibrações excessivas &
\cellcolor{okbg} & \cellcolor{nokbg} & \cellcolor{nabg} \\ \hline
9.2.1.4 & Ventilação adequada ao redor do transformador &
\cellcolor{okbg} & \cellcolor{nokbg} & \cellcolor{nabg} \\ \hline
6.2.8 & Conexões e bornes sem oxidação, aquecimento ou centelhamento &
\cellcolor{okbg} & \cellcolor{nokbg} & \cellcolor{nabg} \\ \hline
6.4.2 & Aterramento do transformador realizado e em bom estado &
\cellcolor{okbg} & \cellcolor{nokbg} & \cellcolor{nabg} \\ \hline

\rowcolor{secbg}
\multicolumn{5}{|l|}{\textbf{CABINE / SALA DE SUBESTAÇÃO}} \\ \hline
9.1.6 & Acesso restrito a pessoal autorizado (cadeado, identificação) &
\cellcolor{okbg} & \cellcolor{nokbg} & \cellcolor{nabg} \\ \hline
9.1.9 & Sinalização de risco elétrico fixada &
\cellcolor{okbg} & \cellcolor{nokbg} & \cellcolor{nabg} \\ \hline
9.2.1.2 & Espaço livre adequado para manobra e manutenção &
\cellcolor{okbg} & \cellcolor{nokbg} & \cellcolor{nabg} \\ \hline
6.2.2 & Eletrodutos instalados de PVC &
\cellcolor{okbg} & \cellcolor{nokbg} & \cellcolor{nabg} \\ \hline
9.2.1.3 & Iluminação adequada no interior da cabine &
\cellcolor{okbg} & \cellcolor{nokbg} & \cellcolor{nabg} \\ \hline
--- & Extintor de incêndio CO$_2$ disponível e dentro do prazo &
\cellcolor{okbg} & \cellcolor{nokbg} & \cellcolor{nabg} \\ \hline
5.2.2 & Eletrodutos vedados com espuma expansiva antichamas &
\cellcolor{okbg} & \cellcolor{nokbg} & \cellcolor{nabg} \\ \hline

\rowcolor{secbg}
\multicolumn{5}{|l|}{\textbf{PROTEÇÕES E SECCIONAMENTO}} \\ \hline
6.3.6 & Chave seccionadora / fusíveis de AT presentes e em bom estado &
\cellcolor{okbg} & \cellcolor{nokbg} & \cellcolor{nabg} \\ \hline
5.4 & Para-raios instalados (proteção contra surto de AT) &
\cellcolor{okbg} & \cellcolor{nokbg} & \cellcolor{nabg} \\ \hline
6.1.5 & Dispositivos de proteção em MT identificados e acessíveis &
\cellcolor{okbg} & \cellcolor{nokbg} & \cellcolor{nabg} \\ \hline

\rowcolor{secbg}
\multicolumn{5}{|l|}{\textbf{ATERRAMENTO E MALHA DE TERRA}} \\ \hline
6.4.2 & Malha de aterramento implantada &
\cellcolor{okbg} & \cellcolor{nokbg} & \cellcolor{nabg} \\ \hline
6.4.4 & Equipotencialização das estruturas metálicas da subestação &
\cellcolor{okbg} & \cellcolor{nokbg} & \cellcolor{nabg} \\ \hline

\end{tabular}%
}}

\vspace{0.3em}
\noindent{\scriptsize \textit{Nota: itens marcados com "---" não são cobertos diretamente pela NBR 14039:2021 — medidor e ramal seguem as normas técnicas da distribuidora (NDU Coelba); extintor de incêndio é regulamentado pelo Corpo de Bombeiros.}}